\chapter*{Reading this edition}
\addcontentsline{toc}{chapter}{Reading this edition}
The Hodge conjecture asks when a rational cohomology class of the correct Hodge type is a rational linear combination of algebraic cycle classes. A useful research account must keep three questions visible. What identifies the target class? Which geometric object supplies a cycle? What proves that the supplied cycle has a nonzero component in the desired rational subspace?

This edition follows that sequence through the projector calculus and the constructive Fermat program, then incorporates the later geometric carriers, signed transfer, global matrix obstructions, explicit standard surfaces, and Weil-class audit. Mathematical developments are retained at theorem and proof level where their stated hypotheses suffice. Proposed uniform arguments are retained as manuscript claims with their outstanding obligations. Failed proposals are part of the evidence.

The August 6 manuscript contained 86 pages. The later August 15 status note downgraded several uniform and correspondence-level claims. The September continuation added a normal weighted carrier at degree 70, a transfer to a full-conductor degree-210 packet, dual certificates excluding a universal carrier method, and explicit recomputations of Aoki's known five-phase cycles. A separate September 22 Weil computation tested a counterexample idea and found replacement algebraic factorizations. This reader reconciles those sources rather than treating the earliest strongest labels as authoritative.

\begin{claimboundary}
No general solution of the rational Hodge conjecture is proved here. The projector equation is an exact reformulation. A finite rank check is a proof of algebraicity only when its vectors are verified algebraic cycles in the specified cohomology space. A balanced Fermat tuple identifies Hodge type and does not construct a cycle. A matching CM sign identifies a possible factor and does not construct a Chow correspondence. An obstructed representative does not establish nonalgebraicity of its class.
\end{claimboundary}

\section*{The mathematical destinations}
For an arbitrary smooth projective variety, the geometric target is a spanning collection of algebraic cycles or an invertible algebraic lowering-and-return transfer. The analytic target is the vanishing of the complete finite integrability obstruction after an allowed stabilization. Neither universal target is obtained here.

For Fermat fourfolds, diagonal automorphisms supply exact residue addresses, rational Galois packets, conductor projectors and finite balance equations. The constructive problem asks for equations, correspondences and a nonzero coefficient that can be followed back to a specified packet. Kang's published result as cited in the checked literature is separated from the explicit algorithms under development. The original proof of that result was not independently re-audited in this edition.

The later carrier reaches a 24-dimensional rational packet at degree 70 at manuscript proof level. The signed identity then reaches a distinct 48-dimensional full-conductor degree-210 packet, using classical Fermat correspondence operations. These positive results coexist with a universal obstruction to the direct positive-anticanonical matrix constructor. That obstruction is a theorem about the method. It is compatible with algebraicity and is bypassed by actual five-phase surfaces.

\section*{Status language}
\begin{longtable}{p{0.23\linewidth}p{0.68\linewidth}}
\toprule
Label & Meaning in this reader\\
\midrule
Exact elementary result & A proof is given using the stated finite-dimensional, integer or polynomial hypotheses.\\
Classical input & An external theorem is invoked at its actual scope. Newly checked leaves are identified in the source notes.\\
Manuscript claim & A geometric or uniform argument is preserved, with its dependency and review boundary. Finite checks do not certify its geometry.\\
Conditional construction & The conclusion follows once the named correspondence, packet containment or coverage hypothesis has been supplied.\\
Finite replay & The implementation was run over the stated inputs and range. No extrapolation beyond that range is made.\\
Failure or retraction & A candidate, inference or universal constructor failed. Its consequence is stated at the level justified by the failure.\\
Open obligation & The required proof object has not been constructed or verified.\\
\bottomrule
\end{longtable}

\section*{A route through the book}
Begin with the rational target, primitive reduction, corrected cycle-frame certificate, support test and transfer theorem. These chapters explain how a geometric construction would become a proof. The analytic branch explains why solving only a projected equation leaves a real gap.

The Fermat chapters separate conductor descent, pairs, surface saturation, Delsarte maps, odd primes and dyadic refinement. The tree derivation is preserved with its later downgrade. The K23 discussion identifies the field and projector typing needed before a rank-one statement can become a rational Chow map.

The current constructive chapters can also be read independently. The degree-70 carrier gives its matrix, hypersurface, singular locus, torus descent, rational Hodge lift and return. The degree-210 identity makes each signed auxiliary factor explicit and shows why actual complementary-pair removal is sufficient. Global capacities have exact finite dual certificates. The cycle chapter verifies containment, a determinant identity, nonzero periods and rational orbit reconstruction for a known standard family. The Weil chapter makes the failed counterexample route concrete.

\section*{The fresh evidence}
The October replay passed the 24 grouped carrier checks, 46 grouped transfer checks and 137 grouped capacity/cycle assertions. It independently reconstructed 364 valid K23 sign certificates and all dyadic census records for moduli 8 through 128. It checked 87,376 weighted dyadic coefficient identities for $3\le s\le9$. The archival 89,478,480-count dyadic report has a larger range and was not replayed in full. The archived projector toy models use floating-point checks with explicit tolerances. The current figures use exact retained data or explicitly schematic structure, and are not themselves algebraic cycle constructions on an actual variety.

The preserved Lean source has 17 theorem scripts and no local \texttt{sorry}, \texttt{admit} or \texttt{axiom} after comments are stripped. Its final theorem assumes an abstract algebraic-packet predicate, seed algebraicity and closure properties. It does not define the needed Chow groups or geometric correspondences. A current build was blocked before elaboration because the restored compiler could not locate its application and the Mathlib checkout was unavailable. This edition records source inspection, rather than a fresh compilation pass.
